\chapter{Sexually Transmitted Infection (STI) Panel}

\section*{What Is This Test?}
A sexually transmitted infection (STI) panel is a bundled set of laboratory tests that simultaneously screens for or diagnoses multiple STIs from a single patient visit. The composition of the panel varies by laboratory, clinical setting, patient demographics, and sexual history. Core tests commonly included are: \textit{Chlamydia trachomatis} and \textit{Neisseria gonorrhoeae} NAAT from genital and extragenital sites (vaginal swab, first-catch urine, rectal swab, pharyngeal swab), syphilis serology (RPR with treponemal-specific confirmatory test or reverse-sequence algorithm using treponemal immunoassay), HIV screening (4th-generation antigen/antibody combination immunoassay), and hepatitis B screening (HBsAg) and hepatitis C screening (anti-HCV with reflex HCV RNA). Expanded panels may include: \textit{Trichomonas vaginalis} NAAT, HSV-2 type-specific serology (in selected patients), and hepatitis A serology (HAV IgG). STI panels are not one-size-fits-all; they must be tailored to the individual's anatomy (specimen sites for NAAT depend on sexual practices---patients reporting receptive anal sex require rectal testing; those reporting receptive oral sex require pharyngeal testing) and individual risk factors.

\section*{Alternative Names}
STD Panel; STI Screening Panel; Comprehensive STI Testing; Sexual Health Panel; STI Workup; GC/Chlamydia/Syphilis/HIV Panel; Full STI Workup

\section*{Principle}
The STI panel is not a single assay but a coordinated battery employing multiple analytic platforms, described in detail in the individual test chapters. Briefly: \textit{C. trachomatis} and \textit{N. gonorrhoeae} NAAT (PCR/TMA) on urine and swab specimens. Syphilis: RPR (nontreponemal flocculation) with reflexive treponemal-specific EIA/CIA and/or TPPA confirmation, using either the traditional or reverse screening algorithm. HIV: 4th-generation chemiluminescent antigen/antibody combination immunoassay with reflex HIV-1/2 antibody differentiation assay and, when indicated, HIV-1 RNA NAT per the CDC testing algorithm. Hepatitis B: HBsAg by CLIA/EIA; if positive, reflex to HBeAg, anti-HBe, and HBV DNA. Hepatitis C: anti-HCV by CLIA/EIA; if positive, reflex HCV RNA quantitative PCR. Additional components as indicated: \textit{T. vaginalis} NAAT, HSV-2 type-specific IgG, HAV total/IgG. Results from the different components are reported separately with turnaround times ranging from 1--2 hours (rapid HIV/GC/CT) to several days (syphilis serology, hepatitis serology).

\section*{History}
Comprehensive STI clinic-based testing evolved from the VD (venereal disease) clinics of the early 20th century, which focused on syphilis and gonorrhea. The development of chlamydia diagnostics (cell culture 1960s, antigen tests 1980s, NAATs 1990s), HIV testing (ELISA 1985, 4th-generation 2000s), and multiplex platforms allowed single-specimen testing for multiple pathogens. The CDC's STI Treatment Guidelines (updated iteratively, most recently 2021) have defined the components of an evidence-based STI evaluation. In the past 5--10 years, extragenital (rectal, pharyngeal) testing has become standard for MSM and others at risk given the high prevalence of asymptomatic extragenital gonorrhea and chlamydia. Self-collected vaginal swabs for GC/CT NAAT have increased screening uptake by eliminating the need for a pelvic speculum examination. The integration of HIV pre-exposure prophylaxis (PrEP) into routine sexual health care (CDC PrEP guidelines 2014, updated 2021) has made comprehensive STI screening at 3-month intervals the standard of care for individuals on PrEP.

\section*{Why This Test Is Ordered}
\begin{itemize}
  \item Comprehensive evaluation of individuals with genitourinary, anorectal, or pharyngeal symptoms suggestive of STIs (discharge, dysuria, ulcers, lesions, proctitis, pharyngitis)
  \item Routine screening of sexually active individuals at increased risk: MSM (at least annually, ideally every 3--6 months depending on risk), individuals with HIV, commercial sex workers, persons with multiple sexual partners, and persons with a new STI diagnosis
  \item Pre-exposure prophylaxis (PrEP) initiation and quarterly monitoring: HIV testing must be documented as negative before PrEP initiation and every 3 months; STI screening (GC/CT at all exposed sites, syphilis) is recommended at each 3-month PrEP visit
  \item Evaluation of individuals who are sexual contacts of a person diagnosed with an STI (contact tracing/partner notification)
  \item Routine prenatal STI screening panel: HIV (opt-out), syphilis (mandatory by law in all states), HBsAg, chlamydia, gonorrhea, and HCV (for women with risk factors)
  \item Assessment of individuals presenting with syndromes that are often STI-related: pelvic inflammatory disease (PID), epididymitis in men under 35, proctitis/proctocolitis, acute arthritis (disseminated gonococcal infection, reactive arthritis), and perianal condyloma
  \item Periodic screening for individuals with substance use disorders (particularly methamphetamine, which is associated with high-risk sexual behaviors)
\end{itemize}

\section*{Primary vs Secondary Test}
The STI panel serves as a primary screening and diagnostic battery for sexually transmitted infections. Its comprehensiveness is designed to detect asymptomatic and minimally symptomatic infections that would be missed by symptom-prompted testing alone. Extragenital screening is primary---failure to test pharyngeal and rectal sites in individuals reporting receptive oral and anal sex results in a falsely negative (incomplete) STI assessment.

\section*{How This Test Is Performed}
Specimen collection varies by the tests included and the patient's anatomy and sexual history. Blood: One SST tube (5--7 mL) and one EDTA tube (for HIV screening and potential nucleic acid testing). Urine: First-catch urine (10--20 mL, patient should not have voided for at least 1 hour). Vaginal swab: Self-collected or clinician-collected flocked swab. Endocervical swab: If speculum examination is performed. Rectal swab: Self-collected or clinician-collected (blind insertion 2--3 cm into anal canal, rotate against mucosa). Pharyngeal swab: Swab tonsillar pillars and posterior pharynx. Multiple swabs can be pooled for testing in some laboratory platforms (reducing cost) or tested individually (preferred for site-specific detection). All specimens must be labeled with the specific anatomic site. Specimens for NAAT are placed in the manufacturer-specific transport medium; specimens for culture (if also ordered) are placed on appropriate transport media or plated at the bedside.

\section*{What Preparation Is Needed}
Patients should avoid urinating for at least 1 hour before first-catch urine collection. Women should avoid douching, vaginal creams, suppositories, and spermicides for 24--48 hours before cervical or vaginal specimen collection. Patients should avoid oral intake for 30 minutes before pharyngeal specimen collection. If possible, patients should refrain from sexual activity for 24--48 hours before specimen collection (especially if rectal testing is performed) to avoid temporary reduction in organism burden. Recent antibiotic use should be reported; testing should ideally be deferred for 3 weeks after completion of antibiotics to avoid false-negative results.

\section*{Contraindications and Precautions}
The STI panel does NOT include testing for all STI pathogens. Notably, \textit{Mycoplasma genitalium} (an emerging cause of persistent/recurrent NGU and cervicitis), HSV-1/HSV-2 (unless specifically ordered), HPV (included only as part of cervical cancer screening per guidelines, not as an STI test), molluscum contagiosum, pubic lice (phthiriasis), scabies, and bacterial vaginosis (not an STI per se but associated with sexual activity) are often excluded from standard panels and must be ordered separately if indicated. A negative STI panel does NOT mean the patient is "STI-free" if only urogenital sites were tested and extragenital exposures were not queried and tested. The window period must be considered: HIV (can be negative up to 3--4 weeks with 4th-generation test), syphilis (RPR may be nonreactive in primary syphilis for the first 1--2 weeks of the chancre). A negative STI panel today does not reflect an exposure that occurred yesterday. The limitations of the tests included must be communicated to the patient.

\section*{Risks}
Venipuncture, swab collection, and urine collection carry minimal to no risk. Psychological distress while awaiting results is common. The risk of a false-positive result exists for all components; the CDC testing algorithms (HIV confirmatory algorithm, syphilis reverse algorithm with TPPA resolution of discordant results) are designed to minimize false-positive reporting. Legal requirements for reporting certain STIs (HIV, syphilis, gonorrhea, chlamydia, chancroid to the health department) must be disclosed.

\section*{What Happens After the Test}
Results are available within hours (rapid HIV, GC/CT NAAT) to days (syphilis serology, hepatitis serology). Reactive/positive results must be communicated to the patient promptly. For notifiable infections (chlamydia, gonorrhea, syphilis, HIV, chancroid), the case must be reported to the local health department. Patients diagnosed with any STI should be screened for other STIs (if not already done). Sexual partners must be evaluated and treated per CDC recommendations. EPT should be offered where legal. Patients treated for bacterial STIs should be re-tested at 3 months (reinfection surveillance). Patients with a reactive HIV test must be linked to HIV care immediately.

\section*{Understanding Results}

\subsection*{Below Normal}
\begin{itemize}
  \item All panel components negative/nonreactive (no STI detected): Indicates no laboratory evidence of the STIs tested for at the tested sites in adequate specimens. Bear in mind the following caveats: (a) a negative HIV test within the window period (0--4 weeks from exposure) does not exclude infection; (b) a nonreactive RPR in a patient with a genital ulcer does not exclude primary syphilis (can be negative in the first 1--2 weeks); (c) negative urogenital GC/CT does not exclude rectal or pharyngeal infection; (d) a negative test panel today does not address exposures from the past 1--4 weeks
\end{itemize}

\subsection*{Above Normal}
\begin{itemize}
  \item Any component positive/reactive: See the individual test chapters for detailed interpretation of each positive result:
    \begin{itemize}
      \item \textit{Chlamydia trachomatis} NAAT positive: Active chlamydial infection at the tested site; treat with doxycycline 100 mg BID x 7 days
      \item \textit{Neisseria gonorrhoeae} NAAT positive: Active gonococcal infection; treat with ceftriaxone 500 mg IM plus doxycycline for chlamydia coverage
      \item Syphilis: RPR reactive with positive treponemal test: Active or past syphilis; stage and treat accordingly
      \item HIV: 4th-generation reactive with positive differentiation assay and/or positive HIV RNA: HIV infection; link to HIV care immediately
      \item HBsAg positive: Active HBV infection (acute or chronic); evaluate HBV DNA, HBeAg/anti-HBe, ALT, and liver fibrosis
      \item Anti-HCV positive with detectable HCV RNA: Active HCV infection; link to care for DAA therapy
    \end{itemize}
  \item Co-infections (2 or more pathogens detected simultaneously): Common, particularly among MSM and individuals with HIV. Each infection should be treated per CDC guidelines; presence of one STI does not alter the treatment of the other, though timing of syphilis staging may be affected if a genital ulcer is explained by HSV
\end{itemize}

\section*{Normal Results}
\textbf{All panel components negative/nonreactive.} No evidence of the STIs tested. The specific results are detailed in each individual test chapter.

\section*{Follow-up Tests}
\begin{itemize}
  \item \textbf{Extragenital (pharyngeal, rectal) GC/CT NAAT if not performed:} Essential for patients reporting receptive oral or anal sex
  \item \textbf{HIV testing (if negative and within window period):} Repeat at 4--6 weeks and 3 months post-exposure
  \item \textbf{Repeat syphilis serology at 4--6 weeks and 3 months:} If negative and high-risk exposure occurred (or if treating empirically for syphilis exposure)
  \item \textbf{Re-testing for chlamydia, gonorrhea, trichomoniasis at 3 months:} Reinfections are common (10--20\%) and warrant routine re-screening, not because of treatment failure
  \item \textit{\textbf{Mycoplasma genitalium}} \textbf{NAAT:} For persistent or recurrent urethritis, cervicitis, or PID after treatment for GC/CT if not part of the initial panel. \textit{M. genitalium} is an emerging pathogen associated with macrolide resistance
  \item \textbf{HSV-2 type-specific serology:} May be added for patients with recurrent genital ulceration or atypical symptoms; not routinely part of STI panels
  \item \textbf{HIV PrEP initiation and continuation monitoring:} For HIV-negative individuals at ongoing high risk, discussion of PrEP with quarterly STI screening + HIV testing + renal function monitoring is standard of care
\end{itemize}

\section*{References}
\begin{enumerate}
  \item Workowski KA, Bachmann LH, Chan PA, et al. Sexually transmitted infections treatment guidelines, 2021. MMWR Recomm Rep. 2021;70(4):1-187.
  \item CDC. Recommendations for the laboratory-based detection of \textit{Chlamydia trachomatis} and \textit{Neisseria gonorrhoeae} --- 2014. MMWR Recomm Rep. 2014;63(RR-02):1-19.
  \item US Preventive Services Task Force. Screening for HIV infection: US Preventive Services Task Force recommendation statement. JAMA. 2019;321(23):2326-2336.
  \item CDC. Preexposure prophylaxis for the prevention of HIV infection in the United States --- 2021 Update: a clinical practice guideline. CDC; 2021.
  \item Tuddenham S, Hamill MM, Ghanem KG. Diagnosis and treatment of sexually transmitted infections: a review. JAMA. 2022;327(2):161-172.
\end{enumerate}

\clearpage
